%
The following blinding procedure is used: in each analysis 
channel, those bins of the MVA (BDT or GNN) discriminant distribution with a signal-to-background ratio ($S/B$) 
greater than 5\% are removed. The procedure applied assuming a given signal at the time to contribute to the corresponding MVA score: for example, when looking at the BDT score evaluated for \mH=500~\gev\ only the presence of a possible signal with \mH=500~\gev\ is accounted for in the blinding procedure. 
The signal is normalised to 12~\ifb that represents the measured four-top-quark cross section from Ref.~\cite{ATLAS-CONF-2021-013} 
minus the expected SM prediction.  

Following this procedure, the number of blinded MVA bins depends marginally on \mH discriminant under consideration; between 8 and 12 bins are blinded overall across all channels for each \mH hypothesis. 
